% ============================================================================
%  ECR methodology - middle-weight version: the five relations you would want
%  to cite are displayed and numbered, everything else runs inline.
%  Follows on from \section{ECR Data-Driven Targeting}\label{sec:ecr-method}.
%  Requires amsmath, amssymb, bm.
%
%  Other variants in this folder:
%    ecr_method_math_2p.tex    two pages, every relation displayed (13 eqs)
%    ecr_method_math_full.tex  six pages, adds the target geometry, the
%                              training algorithm and the remarks on delay
%                              coordinates and deterministic deadlock
%
%  NOTATION: the observed (controller) state is bold, \mathbf{z}_n; plain
%  x, y, z are the Lorenz coordinates.  Rename to \bm{\zeta}_n if the clash
%  with the Lorenz z is unwelcome.
% ============================================================================

\subsection{Plant, section and admissible perturbations}
\label{subsec:ecr-plant}

The Lorenz system is written in full, since the second equation is what defines
the surface of section used throughout,
%
\begin{equation}
\dot{x}=\sigma\,(y-x),
\qquad
\dot{y}=\rho\,x-y-xz,
\qquad
\dot{z}=xy-\beta z,
\label{eq:lorenz}
\end{equation}
%
with state $\mathbf{u}=[x\;y\;z]^{\mathsf T}$, flow
$\varphi_t(\mathbf{u};\mathbf{p})$ and control parameters
$\mathbf{p}=[\sigma\;\rho\;\beta]^{\mathsf T}\in\mathbb{R}^{r}$, $r=3$. The
section is the half-plane
$\Sigma=\{\mathbf{u}:y=y_{\mathrm{PSS}},\,\dot y<0\}$ with
$y_{\mathrm{PSS}}=\sqrt{\beta_{\mathrm{nom}}(\rho_{\mathrm{nom}}-1)}=8.4853$;
transversality requires $\dot y|_{\Sigma}\neq0$ by \eqref{eq:lorenz}, and the
sign condition selects the branch that excludes the equilibrium $C^{+}$. With
the first return time
$\tau=\min\{t>0:\varphi_t(\mathbf{u};\mathbf{p})\in\Sigma\}$ and
$y\equiv y_{\mathrm{PSS}}$ on $\Sigma$, the observed state at the $n$-th
crossing is $\mathbf{z}_n=[x(t_n)\;z(t_n)]^{\mathsf T}\in\mathbb{R}^{N}$,
$N=2$, and the dynamics reduce to the discrete map
%
\begin{equation}
\mathbf{z}_{n+1}=\mathbf{G}(\mathbf{z}_n,\mathbf{p}_n),
\label{eq:map}
\end{equation}
%
the parameter being held constant between crossings. Where only a scalar
measurement $s(t)$ is available,
the delay vector $[s(t_n)\;s(t_n-T)\;\cdots]^{\mathsf T}$ replaces the section
state and \eqref{eq:map} still holds generically.

Each parameter is confined to
$\Pi=\{\mathbf{p}:|p^{i}-p^{i}_{\mathrm{nom}}|<\delta p^{i}_{\max},\,i=1,\dots,r\}$,
here $\mathbf{p}_{\mathrm{nom}}=[10\;28\;8/3]^{\mathsf T}$ and
$\delta\mathbf{p}_{\max}=[0.30\;0.84\;0.08]^{\mathsf T}$, i.e.\ $3\%$ of each
nominal value. That budget is a design choice carried over from the source
paper, not a property of the method: $\Pi$ need only be small enough to leave
the attractor intact, and both the extent of the control regions and the
reaching time scale with how large it is. With the normalised perturbation
$\Delta\mathbf{p}_n=(\mathbf{p}_n-\mathbf{p}_{\mathrm{nom}})\oslash\delta\mathbf{p}_{\max}$,
admissibility is the box test $\|\Delta\mathbf{p}_n\|_{\infty}\le1$. The target
is a fixed point of the controlled map at nominal parameters,
$\mathbf{z}^{*}=\mathbf{G}(\mathbf{z}^{*},\mathbf{p}_{\mathrm{nom}})$ --- a
period-1 unstable periodic orbit of the flow --- located by Newton iteration on
$\mathbf{G}(\mathbf{z},\mathbf{p}_{\mathrm{nom}})-\mathbf{z}$ with a
numerically differentiated Jacobian:
$\mathbf{z}^{*}=[14.2522\;\,39.7867]^{\mathsf T}$, with map eigenvalues
$\lambda_1=4.71$ and $\lambda_2\approx10^{-9}$, so the return map is
effectively one-dimensional. The displacement reachable within $\Pi$ exceeds
the target radius $\delta=0.30$ tenfold, which makes single-step targeting
feasible.

\subsection{Control regions and their empirical construction}
\label{subsec:ecr-regions}

With $\mathcal{B}_\delta=\{\mathbf{z}:\|\mathbf{z}-\mathbf{z}^{*}\|<\delta\}$,
the target region and the extended control regions are
%
\begin{align}
T_0&=\Bigl\{\mathbf{z}\in\mathcal{B}_\delta:\exists\,\mathbf{p}\in\Pi,\;
\mathbf{G}(\mathbf{z},\mathbf{p})\in\mathcal{B}_\delta\Bigr\},
\label{eq:def-T0}\\[2pt]
T_i&=\Bigl\{\mathbf{z}\notin\textstyle\bigcup_{j<i}T_j:\exists\,\mathbf{p}\in\Pi,\;
\mathbf{G}(\mathbf{z},\mathbf{p})\in T_{i-1}\Bigr\},
\qquad i=1,\dots,K .
\label{eq:def-Ti}
\end{align}
%
The regions are disjoint by construction, membership of $T_i$ certifies capture
of $T_0$ within $i$ admissible steps, and their union is the activation region,
outside which $\mathbf{p}_{\mathrm{nom}}$ is applied. Both are existence
statements, settled empirically: the plant is run from many initial conditions
with an independent $\mathbf{p}_n\sim\mathcal{U}(\Pi)$ applied at each
crossing, giving the transitions
$\mathcal{D}=\{(\mathbf{z}_n,\mathbf{p}_n,\mathbf{z}_{n+1})\}_{n=1}^{M}$ in
which the applied $\mathbf{p}_n$ is the witness. Level zero follows directly
from \eqref{eq:def-T0} as
$\mathcal{I}_0=\{n:\|\mathbf{z}_n-\mathbf{z}^{*}\|<\delta\wedge
\|\mathbf{z}_{n+1}-\mathbf{z}^{*}\|<\delta\}$, and for $i\ge1$
$\mathcal{I}_i=\{n:\mathcal{M}_{i-1}(\mathbf{z}_{n+1})=1\wedge
\mathcal{M}_{j}(\mathbf{z}_{n})=0\;\forall j<i\}$, membership of the successor
in $T_{i-1}$ being decided by the already-trained level-$(i-1)$ model acting as
an oracle $\mathcal{M}_{i-1}:\mathbb{R}^{N}\to\{0,1\}$, defined below. Region
$i$ is thus trained on
$\{(\mathbf{z}_n,\Delta\mathbf{p}_n)\}_{n\in\mathcal{I}_i}$ --- the state, and
the perturbation observed to move it one region inward --- bottom-up, each
level supplying the oracle for the next.

\subsection{Clustering, region assignment and local models}
\label{subsec:ecr-clustering}

ECR-I fits one network over the whole activation region and ECR-II one per
region; ECR-III partitions each region first, since data in the outer regions
is scattered. On $\mathcal{Z}_i=\{\mathbf{z}_n\}_{n\in\mathcal{I}_i}$ the clusters
$C_{ij}$ are the equivalence classes of the transitive closure of
$\mathbf{z}\sim\mathbf{z}'\iff\|\mathbf{z}-\mathbf{z}'\|\le r_i$, with $r_i$
taken at the first minimum of the histogram of pairwise distances. Each cluster
carries a mean $\bm{\mu}_{ij}$ and a ridge-regularised covariance
$\bm{\Sigma}_{ij}$, and a state is assigned to the cluster minimising the
normalised Mahalanobis distance
%
\begin{equation}
\mathrm{NMD}^{2}_{ij}(\mathbf{z}_n)=(\mathbf{z}_n-\bm{\mu}_{ij})^{\mathsf T}\,
\underline{\mathbf{N}}_{ij}\,(\mathbf{z}_n-\bm{\mu}_{ij}),
\qquad
\underline{\mathbf{N}}_{ij}=\Bigl(\bm{\Sigma}_{ij}\big/\det(\bm{\Sigma}_{ij})^{1/N}\Bigr)^{-1},
\label{eq:nmd}
\end{equation}
%
the unit-determinant normalisation removing the size of a cluster but not its
shape, so clusters of different populations compete on equal terms (plain,
variance- and trace-normalised readings are alternatives). Each cluster
also carries a Gaussian RBF network mapping the standardised state
$\tilde{\mathbf{z}}$ to a normalised perturbation,
$\widehat{\Delta\mathbf{p}}(\mathbf{z})=\mathbf{W}^{\mathsf T}[\phi_1\;\cdots\;\phi_{n_c}\;1]^{\mathsf T}$
with $\phi_k=\exp(-\|\tilde{\mathbf{z}}-\mathbf{c}_k\|^{2}/2s_k^{2})$, centres
from $k$-means and weights
$\mathbf{W}=(\bm{\Phi}^{\mathsf T}\bm{\Phi}+\lambda\mathbf{I})^{-1}\bm{\Phi}^{\mathsf T}\mathbf{Y}$
in closed form. The oracle used above is an admissibility test on the nearest
cluster's network,
$\mathcal{M}_i(\mathbf{z})=\mathbb{1}[\|\widehat{\Delta\mathbf{p}}^{\,(ij^{*})}(\mathbf{z})\|_{\infty}\le1]$
with $j^{*}=\arg\min_j\mathrm{NMD}_{ij}(\mathbf{z})$, optionally gated on the
cluster's own training NMD, since \eqref{eq:nmd} returns a nearest cluster even
for states far outside the data.

\subsection{Control law and performance measures}
\label{subsec:ecr-online}

At each crossing the state is measured, a region is identified and the
corresponding perturbation applied,
%
\begin{equation}
\mathbf{p}_n=
\begin{cases}
\mathbf{p}_{\mathrm{nom}}+\widehat{\Delta\mathbf{p}}^{\,(i^{*}j^{*})}(\mathbf{z}_n)\odot\delta\mathbf{p}_{\max},
& \text{if an admissible } (i^{*},j^{*}) \text{ exists},\\
\mathbf{p}_{\mathrm{nom}}, & \text{otherwise},
\end{cases}
\label{eq:control-law}
\end{equation}
%
so the applied perturbation is admissible by construction. The variants differ
only in how the pair is found: ECR-II takes the first admissible answer in
ascending region order,
$i^{*}=\min\{i:\|\widehat{\Delta\mathbf{p}}^{\,(i)}(\mathbf{z}_n)\|_{\infty}\le1\}$,
whereas ECR-III obtains it analytically as
$(i^{*},j^{*})=\arg\min_{i,j}\mathrm{NMD}_{ij}(\mathbf{z}_n)$ over the clusters
of all regions at once, trading a sequential search over networks for a
distance computation. Performance is measured by the reaching time and its mean
over initial conditions on the attractor,
%
\begin{equation}
n_r(\mathbf{z}_0)=\min\bigl\{n\ge0:\|\mathbf{z}_n-\mathbf{z}^{*}\|<\delta\bigr\},
\qquad
\bar{n}_r=\frac{1}{M_{\mathrm{ic}}}\sum_{m=1}^{M_{\mathrm{ic}}}
\min\bigl\{n_r(\mathbf{z}_0^{(m)}),n_{\max}\bigr\},
\label{eq:reaching-time}
\end{equation}
%
censored at a horizon $n_{\max}$, together with the capture rate
$\Pr\{n_r<n_{\max}\}$ and the retention --- the fraction of steps after capture
spent inside $\mathcal{B}_\delta$ --- which separates targeting from
stabilisation. Without targeting $\bar{n}_r$ is the natural recurrence time to
$\mathcal{B}_\delta$, some $15$ steps here.
